\documentclass[11pt,paper=a4,answers]{exam}
\usepackage{graphicx,lastpage}
\usepackage{upgreek}
\usepackage{censor}
\usepackage{gensymb}
\usepackage{amsmath}
\usepackage{multicol}
\usepackage{enumitem}
\usepackage{tasks}
\usepackage{adjustbox}
\usepackage{xcolor}
\usepackage{tfrupee}
\setlength{\voffset}{0.25in}
\usepackage{tabularx}
\newcolumntype{Y}{>{\centering\arraybackslash}X}
%==============================================================
% Customise Non-Essential Packages Below
% =============================================================
\usepackage[most]{tcolorbox}
\usepackage{eso-pic}
\usepackage{tikz}
\AddToShipoutPictureBG{%
\begin{tikzpicture}[remember picture,overlay]
\foreach \x in {-8,-4,0,4,8,12,16,20}
\foreach \y in {-8,-4,0,4,8,12,16,20} {
\node[opacity=0.05,rotate=45,anchor=center] at ([xshift=\x cm,yshift=\y cm]current page.center) {%
\begin{minipage}{2cm}
\centering
\includegraphics[width=2cm]{images/logo_black.png}\\
\small preppeo.com
\end{minipage}%
};
}
\end{tikzpicture}%
}
\printanswers

%==============================================================
% Customise Non-Essential Packages Above
% =============================================================
\definecolor{svgray}{rgb}{0.90, 0.90, 0.90}
\graphicspath{{images/}}

\usepackage[normalem]{ulem}
\renewcommand\ULthickness{1pt}   %%---> For changing thickness of underline
\setlength\ULdepth{3.5ex}%\maxdimen ---> For changing depth of underline
\renewcommand{\baselinestretch}{1}
\chead{
\includegraphics[scale=0.1,height=2cm ]{logo.png}
}
\pagestyle{empty}

\pagestyle{headandfoot}
\headrule
\cfoot{W: https://preppeo.com; M: +91-8130711689}

\begin{document}
%==============================================================
% Customise Question paper details below this
% =============================================================
\begin{center}
    \centering
    \renewcommand{\arraystretch}{1.5}
    \begin{tabularx}{\textwidth}{YYY}
    & \normalsize\bf{{\Large{{Quantitative Reasoning}}}} & \\
    By Preppeo & \normalsize\bf{{sat\_lid\_028}} & SAT\\
    \normalsize{{\textbf{{Unit:}} Advanced Math}} & \normalsize{{\textbf{{Chapter:}} Function Notation	Function notation and interpreting functions}} & \normalsize{{\textbf{{Topic:}} Function Transformations}} \\
    \end{tabularx}
\end{center}
\par\noindent\rule{\textwidth}{1pt}
% =============================================================
% Customise Question paper details above this
% =============================================================

\begin{questions}
\pointsinrightmargin
\pointsdroppedatright
\marksnotpoints
\pointpoints{mark}{marks}
\pointformat{\boldmath\themarginpoints}
\bracketedpoints

% =============================================================
% Question Begin
% =============================================================

% Lid Topic: Advanced Math — Function Transformations
% Total Questions: 50

% --- PART 1: Translations and Reflections (1-25) ---

% Question 1
% Category: Concept | Difficulty: Medium | Question Type: Multiple Choice
\question
The graph of $y = f(x)$ is shown in the $xy$-plane. Which of the following equations represents a graph that is the result of shifting the graph of $f$ upward by 4 units?
\begin{tasks}(2)
    \task $y = f(x + 4)$
    \task $y = f(x - 4)$
    \task $y = f(x) + 4$
    \task $y = f(x) - 4$
\end{tasks}

\begin{solution}
\textbf{Conceptual Explanation:} \\
A vertical shift is achieved by adding or subtracting a constant outside the function notation. Adding a positive constant $k$ shifts the graph upward.

\vspace{25pt}

Answer: \colorbox{yellow!100}{C}
\end{solution}

\vspace{40pt}

% Question 2
% Category: Exam level | Difficulty: Hard | Question Type: Numeric Entry
\question
\begin{center}
\begin{tikzpicture}[scale=0.6]
    \draw[->] (-3,0) -- (3,0) node[right] {$x$};
    \draw[->] (0,-1) -- (0,5) node[above] {$y$};
    \draw[domain=-2:2, smooth, variable=\x, blue, thick] plot ({\x}, {\x*\x});
    \filldraw[black] (0,0) circle (2pt) node[below] {$(0, 0)$};
    \node at (1.5,4) {$f(x)$};
\end{tikzpicture}
\end{center}
The vertex of the parabola $f(x) = x^2$ is $(0, 0)$. If $g(x) = f(x-5) + 3$, what is the $y$-coordinate of the vertex of the graph of $g$?

\begin{solution}
\textbf{Conceptual Explanation:} \\
The transformation $f(x-h) + k$ shifts the original graph $h$ units horizontally and $k$ units vertically. 

\vspace{20pt}

\textbf{Calculation and Logic:} \\
Original vertex: $(0, 0)$. \\
Horizontal shift: $x-5$ means 5 units to the right ($h=5$). \\
Vertical shift: $+3$ means 3 units up ($k=3$). \\
New vertex: $(5, 3)$.

\vspace{10pt}

The $y$-coordinate is 3.

\vspace{25pt}

Answer: \colorbox{yellow!100}{3}
\end{solution}

\vspace{40pt}

% Question 3
% Category: Practice | Difficulty: Medium | Question Type: Multiple Choice
\question
If the point $(2, 10)$ is on the graph of $y = f(x)$, which point must be on the graph of $y = f(x + 3)$?
\begin{tasks}(4)
    \task $(5, 10)$
    \task $(-1, 10)$
    \task $(2, 13)$
    \task $(2, 7)$
\end{tasks}

\begin{solution}
\textbf{Conceptual Explanation:} \\
The transformation $f(x + h)$ represents a horizontal shift to the left by $h$ units. This affects only the $x$-coordinate by subtracting $h$.

\vspace{20pt}

\textbf{Calculation and Logic:} \\
Original $x = 2$. \\
Shift left 3 units: $2 - 3 = -1$. \\
The $y$-coordinate remains 10. \\
New point: $(-1, 10)$.

\vspace{25pt}

Answer: \colorbox{yellow!100}{B}
\end{solution}

\vspace{40pt}

% Question 4
% Category: Exam level | Difficulty: Hard | Question Type: Numeric Entry
\question
The function $g$ is defined by $g(x) = f(x) - 7$. If the $y$-intercept of the graph of $f$ is $(0, 15)$, what is the $y$-coordinate of the $y$-intercept of the graph of $g$?

\begin{solution}
\textbf{Conceptual Explanation:} \\
A vertical shift affects the $y$-intercept directly. If a graph moves down 7 units, its $y$-intercept also moves down 7 units.

\vspace{20pt}

\textbf{Calculation and Logic:} \\
$g(0) = f(0) - 7$. \\
Since $f(0) = 15$: \\
$g(0) = 15 - 7 = 8$.

\vspace{25pt}

Answer: \colorbox{yellow!100}{8}
\end{solution}

\vspace{40pt}

% Question 5
% Category: Concept | Difficulty: Medium | Question Type: Multiple Choice
\question
\begin{center}
\begin{tikzpicture}[scale=0.6]
    \draw[->] (-1,0) -- (5,0) node[right] {$x$};
    \draw[->] (0,-1) -- (0,5) node[above] {$y$};
    \draw[domain=0.5:4.5, smooth, variable=\x, red, thick] plot ({\x}, {sqrt(\x)});
    \node[red] at (4,1.5) {$f$};
\end{tikzpicture}
\end{center}
Which transformation of function $f$ would result in a reflection across the $x$-axis?
\begin{tasks}(2)
    \task $y = f(-x)$
    \task $y = -f(x)$
    \task $y = f(x) - 1$
    \task $y = \frac{1}{f(x)}$
\end{tasks}

\begin{solution}
\textbf{Conceptual Explanation:} \\
Reflecting a graph across the $x$-axis means negating all of its $y$-values. This is represented by placing a negative sign outside the function: $-f(x)$.

\vspace{25pt}

Answer: \colorbox{yellow!100}{B}
\end{solution}

\vspace{40pt}

% Question 6
% Category: Exam level | Difficulty: Hard | Question Type: Numeric Entry
\question
A line has a $y$-intercept of $(0, 4)$ and a slope of 2. The line is shifted 3 units to the right to create a new line. What is the $y$-intercept of the new line?

\begin{solution}
\textbf{Calculation and Logic:} \\
Original line: $f(x) = 2x + 4$. \\
Shift right 3: $g(x) = f(x - 3) = 2(x - 3) + 4$. \\
Simplify: $g(x) = 2x - 6 + 4 = 2x - 2$. \\
The $y$-intercept of the new line is $-2$.

\vspace{25pt}

Answer: \colorbox{yellow!100}{-2}
\end{solution}

\vspace{40pt}

% Question 7
% Category: Practice | Difficulty: Medium | Question Type: Multiple Choice
\question
$g(x) = f(x + 2) - 5$ \\
The graph of function $g$ is the result of which sequence of transformations on function $f$?
\begin{tasks}(1)
    \task Shift 2 units left and 5 units down
    \task Shift 2 units right and 5 units down
    \task Shift 2 units left and 5 units up
    \task Shift 2 units right and 5 units up
\end{tasks}

\begin{solution}
\textbf{Conceptual Explanation:} \\
$(x+2)$ inside the function shifts left. $-5$ outside shifts down.

\vspace{25pt}

Answer: \colorbox{yellow!100}{A}
\end{solution}

\vspace{40pt}

% Question 8
% Category: Exam level | Difficulty: Hard | Question Type: Numeric Entry
\question
If the point $(a, b)$ is on the graph of $y = f(x)$, and the point $(a+4, b-6)$ is on the graph of $y = f(x-h) + k$, what is the value of $h + k$?

\begin{solution}
\textbf{Calculation and Logic:} \\
Moving from $a$ to $a+4$ is a shift 4 units to the right, so $h = 4$. \\
Moving from $b$ to $b-6$ is a shift 6 units down, so $k = -6$. \\
$h + k = 4 + (-6) = -2$.

\vspace{25pt}

Answer: \colorbox{yellow!100}{-2}
\end{solution}

\vspace{40pt}

% Question 9
% Category: Concept | Difficulty: Medium | Question Type: Multiple Choice
\question

If the graph of $f(x) = x^2$ is reflected across the $x$-axis and then shifted up 10 units, which of the following defines the new function $g$?
\begin{tasks}(2)
    \task $g(x) = -x^2 + 10$
    \task $g(x) = -(x^2 + 10)$
    \task $g(x) = (-x)^2 + 10$
    \task $g(x) = -x^2 - 10$
\end{tasks}

\begin{solution}
\textbf{Calculation and Logic:} \\
Reflect across x-axis: $-f(x) = -x^2$. \\
Shift up 10: $-x^2 + 10$.

\vspace{25pt}

Answer: \colorbox{yellow!100}{A}
\end{solution}

\vspace{40pt}

% Question 10
% Category: Exam level | Difficulty: Hard | Question Type: Numeric Entry
\question
The vertex of a parabola is $(3, -4)$. After a transformation, the new vertex is $(1, 2)$. If the transformation is written as $f(x-h)+k$, what is the value of $h$?

\begin{solution}
\textbf{Calculation and Logic:} \\
Original $x = 3$. New $x = 1$. \\
Shift left 2 units: $3 - 2 = 1$. \\
A shift left 2 is represented as $f(x+2)$, so $h = -2$.

\vspace{25pt}

Answer: \colorbox{yellow!100}{-2}
\end{solution}

\vspace{40pt}

% Question 11
% Category: Practice | Difficulty: Medium | Question Type: Multiple Choice
\question
\begin{center}
\begin{tikzpicture}[scale=0.5]
    \draw[->] (-3,0) -- (3,0) node[right] {$x$};
    \draw[->] (0,-1) -- (0,5) node[above] {$y$};
    \draw[domain=-2:2, smooth, variable=\x, blue, thick] plot ({\x}, {\x*\x + 1});
    \node at (1.5,4.5) {$y=f(x)$};
\end{tikzpicture}
\end{center}
The graph of $y=f(x)$ is shown. Which of the following is the graph of $y=f(x)-3$?
\begin{tasks}(1)
    \task The graph of $f$ shifted 3 units left.
    \task The graph of $f$ shifted 3 units right.
    \task The graph of $f$ shifted 3 units up.
    \task The graph of $f$ shifted 3 units down.
\end{tasks}

\begin{solution}
\textbf{Conceptual Explanation:} \\
Subtracting 3 from the function value moves every point on the graph downward by 3 units.

\vspace{25pt}

Answer: \colorbox{yellow!100}{D}
\end{solution}

\vspace{40pt}

% Question 12
% Category: Exam level | Difficulty: Hard | Question Type: Numeric Entry
\question
$f(x) = 2x - 8$ \\
$g(x) = f(x+k)$ \\
If the $x$-intercept of $g$ is $(0, 0)$, what is the value of $k$?

\begin{solution}
\textbf{Calculation and Logic:} \\
$x$-intercept of $g$ is $(0,0) \Rightarrow g(0) = 0$. \\
$f(0+k) = 0 \Rightarrow f(k) = 0$. \\
$2k - 8 = 0 \Rightarrow 2k = 8 \Rightarrow k = 4$.

\vspace{25pt}

Answer: \colorbox{yellow!100}{4}
\end{solution}

\vspace{40pt}

% Question 13
% Category: Practice | Difficulty: Medium | Question Type: Multiple Choice
\question
How does the graph of $y = 5f(x)$ compare to the graph of $y = f(x)$?
\begin{tasks}(1)
    \task It is shifted 5 units up.
    \task It is shifted 5 units right.
    \task It is vertically stretched by a factor of 5.
    \task It is horizontally compressed by a factor of 5.
\end{tasks}

\begin{solution}
\textbf{Conceptual Explanation:} \\
Multiplying the output of a function by a constant $a > 1$ results in a vertical stretch.

\vspace{25pt}

Answer: \colorbox{yellow!100}{C}
\end{solution}

\vspace{40pt}

% Question 14
% Category: Exam level | Difficulty: Hard | Question Type: Numeric Entry
\question
A function $f$ has a domain of $0 \leq x \leq 10$. What is the width of the domain of the function $g(x) = f(2x)$?

\begin{solution}
\textbf{Conceptual Explanation:} \\
$f(ax)$ results in a horizontal compression if $a > 1$. The $x$-values are divided by $a$.

\vspace{20pt}

\textbf{Calculation and Logic:} \\
Original domain: $[0, 10]$, width = 10. \\
New domain: $0 \leq 2x \leq 10 \Rightarrow 0 \leq x \leq 5$. \\
Width of new domain = 5.

\vspace{25pt}

Answer: \colorbox{yellow!100}{5}
\end{solution}

\vspace{40pt}

% Question 15
% Category: Concept | Difficulty: Medium | Question Type: Multiple Choice
\question
\begin{center}
\begin{tikzpicture}[scale=0.6]
    \draw[->] (-3,0) -- (3,0) node[right] {$x$};
    \draw[->] (0,-3) -- (0,3) node[above] {$y$};
    \draw[thick, blue] (-2,-1) -- (2,2);
    \node at (1.5,1) {$f$};
    \draw[thick, red, dashed] (2,-1) -- (-2,2);
    \node at (-1.5,1) {$g$};
\end{tikzpicture}
\end{center}
Which transformation maps the solid line $f$ to the dashed line $g$?
\begin{tasks}(2)
    \task $g(x) = f(-x)$
    \task $g(x) = -f(x)$
    \task $g(x) = f(x-2)$
    \task $g(x) = f(x+2)$
\end{tasks}

\begin{solution}
\textbf{Conceptual Explanation:} \\
Negating the input $x \rightarrow -x$ reflects the graph across the $y$-axis.

\vspace{25pt}

Answer: \colorbox{yellow!100}{A}
\end{solution}

\vspace{40pt}

% Question 16
% Category: Exam level | Difficulty: Hard | Question Type: Numeric Entry
\question
If the point $(4, 8)$ is on the graph of $y = f(x)$, what is the $y$-coordinate of the corresponding point on the graph of $y = \frac{1}{2}f(x) + 5$?

\begin{solution}
\textbf{Calculation and Logic:} \\
Apply transformations to the $y$-value: \\
1. Vertical compression: $8 \times (1/2) = 4$. \\
2. Vertical shift up: $4 + 5 = 9$.

\vspace{25pt}

Answer: \colorbox{yellow!100}{9}
\end{solution}

\vspace{40pt}

% Question 17
% Category: Practice | Difficulty: Medium | Question Type: Multiple Choice
\question
Which of the following describes the graph of $y = f(x-10)$?
\begin{tasks}(1)
    \task Shift 10 units up
    \task Shift 10 units down
    \task Shift 10 units left
    \task Shift 10 units right
\end{tasks}

\begin{solution}
\textbf{Conceptual Explanation:} \\
Subtracting from the input $x$ shifts the graph horizontally to the right.

\vspace{25pt}

Answer: \colorbox{yellow!100}{D}
\end{solution}

\vspace{40pt}

% Question 18
% Category: Exam level | Difficulty: Hard | Question Type: Numeric Entry
\question
The function $f(x) = |x|$ is shifted 4 units left and 2 units up. What is the value of the new function at $x = 0$?

\begin{solution}
\textbf{Calculation and Logic:} \\
New function: $g(x) = |x + 4| + 2$. \\
Evaluate at $x = 0$: \\
$g(0) = |0 + 4| + 2 = 4 + 2 = 6$.

\vspace{25pt}

Answer: \colorbox{yellow!100}{6}
\end{solution}

\vspace{40pt}

% Question 19
% Category: Concept | Difficulty: Medium | Question Type: Multiple Choice
\question

If the point $(3, 5)$ is on the graph of $y = f(x)$, which point must be on the graph of $y = f(-x)$?
\begin{tasks}(4)
    \task $(3, -5)$
    \task $(-3, 5)$
    \task $(-3, -5)$
    \task $(5, 3)$
\end{tasks}

\begin{solution}
\textbf{Conceptual Explanation:} \\
$f(-x)$ reflects the $x$-coordinates.

\vspace{25pt}

Answer: \colorbox{yellow!100}{B}
\end{solution}

\vspace{40pt}

% Question 20
% Category: Exam level | Difficulty: Hard | Question Type: Numeric Entry
\question
$f(x) = (x-2)^2 + 1$ \\
$g(x) = f(x+3) - 4$ \\
What is the $x$-coordinate of the vertex of $g$?

\begin{solution}
\textbf{Calculation and Logic:} \\
Vertex of $f$: $(2, 1)$. \\
Transformation $g$: shift left 3, down 4. \\
$x$-coordinate: $2 - 3 = -1$.

\vspace{25pt}

Answer: \colorbox{yellow!100}{-1}
\end{solution}

\vspace{40pt}

% Question 21
% Category: Practice | Difficulty: Medium | Question Type: Multiple Choice
\question
How does the graph of $y = f(3x)$ compare to $y = f(x)$?
\begin{tasks}(1)
    \task Vertically stretched by 3
    \task Vertically compressed by 1/3
    \task Horizontally stretched by 3
    \task Horizontally compressed by 1/3
\end{tasks}

\begin{solution}
\textbf{Conceptual Explanation:} \\
Input multipliers $> 1$ compress the graph horizontally.

\vspace{25pt}

Answer: \colorbox{yellow!100}{D}
\end{solution}

\vspace{40pt}

% Question 22
% Category: Exam level | Difficulty: Hard | Question Type: Numeric Entry
\question
If $f(x) = x^2 + 4$, what is the minimum value of $g(x) = f(x-10) + 5$?

\begin{solution}
\textbf{Calculation and Logic:} \\
Minimum of $f$ is 4 (at $x=0$). \\
$g(x)$ is $f$ shifted 10 right and 5 up. \\
The vertical shift up 5 changes the minimum from 4 to $4 + 5 = 9$.

\vspace{25pt}

Answer: \colorbox{yellow!100}{9}
\end{solution}

\vspace{40pt}

% Question 23
% Category: Concept | Difficulty: Medium | Question Type: Multiple Choice
\question
If $g(x) = f(x) + k$, and the graph of $g$ is lower than the graph of $f$, what must be true about $k$?
\begin{tasks}(4)
    \task $k > 0$
    \task $k < 0$
    \task $k = 0$
    \task $k = x$
\end{tasks}

\begin{solution}
\textbf{Conceptual Explanation:} \\
A downward shift occurs when the constant being added is negative.

\vspace{25pt}

Answer: \colorbox{yellow!100}{B}
\end{solution}

\vspace{40pt}

% Question 24
% Category: Exam level | Difficulty: Hard | Question Type: Numeric Entry
\question
If $(2, 4)$ is on the graph of $f(x)$, and $(2, -4)$ is on the graph of $a \cdot f(x)$, what is the value of $a$?

\begin{solution}
\textbf{Calculation and Logic:} \\
$a \cdot f(2) = -4 \Rightarrow a(4) = -4 \Rightarrow a = -1$.

\vspace{25pt}

Answer: \colorbox{yellow!100}{-1}
\end{solution}

\vspace{40pt}

% Question 25
% Category: Practice | Difficulty: Medium | Question Type: Multiple Choice
\question
Which transformation results in a graph that is wider than the original $y = |x|$?
\begin{tasks}(2)
    \task $y = 2|x|$
    \task $y = 0.5|x|$
    \task $y = |x| + 2$
    \task $y = |x+2|$
\end{tasks}

\begin{solution}
\textbf{Conceptual Explanation:} \\
A vertical compression ($0 < a < 1$) makes a graph appear wider.

\vspace{25pt}

Answer: \colorbox{yellow!100}{B}
\end{solution}

% --- PART 2: Advanced Scaling & Multi-Step Transformations (26-50) ---

% Question 26
% Category: Exam level | Difficulty: Hard | Question Type: Numeric Entry
\question
\begin{center}
\begin{tikzpicture}[scale=0.6]
    \draw[->] (-1,0) -- (6,0) node[right] {$x$};
    \draw[->] (0,-1) -- (0,5) node[above] {$y$};
    \draw[domain=0:4.5, smooth, variable=\x, blue, thick] plot ({\x}, {sqrt(\x)});
    \filldraw[black] (4,2) circle (2pt) node[above] {$(4, 2)$};
    \node[blue] at (2,1.8) {$y = f(x)$};
\end{tikzpicture}
\end{center}
The point $(4, 2)$ is on the graph of the function $f$ shown above. If $g(x) = f(x+5) + 6$, what is the $y$-coordinate of the corresponding point on the graph of $g$?

\begin{solution}
\textbf{Conceptual Explanation:} \\
A function transformation in the form $g(x) = f(x-h) + k$ shifts the original coordinates. The value inside the parentheses ($h$) affects the $x$-coordinate, while the value outside ($k$) affects the $y$-coordinate.

\vspace{20pt}

\textbf{Calculation and Logic:} \\
Original point: $(4, 2)$. \\
1. \textbf{Identify the Vertical Transformation:} The $+6$ outside the function indicates a vertical shift upward by 6 units. \\
2. \textbf{Apply to the y-coordinate:} Original $y = 2$. New $y = 2 + 6 = 8$.

\vspace{10pt}

Note: The $(x+5)$ would shift the $x$-coordinate 5 units to the left ($4 - 5 = -1$), but the question specifically asks for the $y$-coordinate.

\vspace{25pt}

Answer: \colorbox{yellow!100}{8}
\end{solution}

\vspace{40pt}

% Question 27
% Category: Practice | Difficulty: Medium | Question Type: Multiple Choice
\question
If $f(x) = x^2$, which of the following defines a function $g$ whose graph is the result of reflecting the graph of $f$ across the $y$-axis?
\begin{tasks}(4)
    \task $g(x) = -x^2$
    \task $g(x) = (-x)^2$
    \task $g(x) = x^2 - 1$
    \task $g(x) = \frac{1}{x^2}$
\end{tasks}

\begin{solution}
\textbf{Conceptual Explanation:} \\
A reflection across the $y$-axis is achieved by replacing the input $x$ with its negative counterpart, $-x$. This is represented as $f(-x)$.

\vspace{20pt}

\textbf{Calculation and Logic:} \\
Given $f(x) = x^2$. \\
Applying the transformation: $g(x) = f(-x) = (-x)^2$. \\
(Interestingly, for $x^2$, the graph remains identical because $(-x)^2 = x^2$, illustrating that parabolas centered at the $y$-axis have $y$-axis symmetry).

\vspace{25pt}

Answer: \colorbox{yellow!100}{B}
\end{solution}

\vspace{40pt}

% Question 28
% Category: Exam level | Difficulty: Hard | Question Type: Numeric Entry
\question
The function $f$ is defined by $f(x) = 3x + 4$. The graph of $g$ is the result of shifting the graph of $f$ down by 10 units. What is the $x$-intercept of the graph of $g$?

\begin{solution}
\textbf{Conceptual Explanation:} \\
First, determine the equation of the new function $g$ by applying the vertical shift. Then, solve for the $x$-intercept by setting $g(x) = 0$.

\vspace{20pt}

\textbf{Calculation and Logic:} \\
1. \textbf{Define g(x):} $g(x) = f(x) - 10 = (3x + 4) - 10 = 3x - 6$. \\
2. \textbf{Find x-intercept:} Set $y = 0$: \\
$0 = 3x - 6$ \\
$6 = 3x$ \\
$x = 2$

\vspace{25pt}

Answer: \colorbox{yellow!100}{2}
\end{solution}

\vspace{40pt}

% Question 29
% Category: Concept | Difficulty: Medium | Question Type: Multiple Choice
\question
\begin{center}
\begin{tikzpicture}[scale=0.5]
    \draw[->] (-4,0) -- (4,0) node[right] {$x$};
    \draw[->] (0,-1) -- (0,5) node[above] {$y$};
    \draw[domain=-2:2, smooth, variable=\x, blue, thick] plot ({\x}, {\x*\x});
    \draw[domain=-1:1, smooth, variable=\x, red, dashed, ultra thick] plot ({\x}, {4*\x*\x});
    \node[blue] at (2.2,4) {$f$};
    \node[red] at (0.5,3.5) {$g$};
\end{tikzpicture}
\end{center}
The graph of $f(x) = x^2$ is transformed into the graph of $g$. Which of the following could define $g$?
\begin{tasks}(2)
    \task $g(x) = \frac{1}{4}x^2$
    \task $g(x) = 4x^2$
    \task $g(x) = x^2 + 4$
    \task $g(x) = (x+4)^2$
\end{tasks}

\begin{solution}
\textbf{Conceptual Explanation:} \\
A graph that becomes "steeper" or "narrower" is undergoing a vertical stretch. This happens when the function is multiplied by a constant $a > 1$.

\vspace{20pt}

\textbf{Calculation and Logic:} \\
Looking at the diagram, the dashed line $g$ is narrower than the original $f$. This indicates the $y$-values are increasing faster for the same $x$-values. Only $g(x) = 4x^2$ represents a vertical stretch.

\vspace{25pt}

Answer: \colorbox{yellow!100}{B}
\end{solution}

\vspace{40pt}

% Question 30
% Category: Exam level | Difficulty: Hard | Question Type: Numeric Entry
\question
If the graph of $y = f(x)$ is shifted 7 units to the right and 2 units up, the resulting graph is $y = f(x-h) + k$. What is the value of $hk$?

\begin{solution}
\textbf{Calculation and Logic:} \\
1. \textbf{Identify h:} A shift 7 units right is represented by $(x - 7)$. Thus, $h = 7$. \\
2. \textbf{Identify k:} A shift 2 units up is represented by $+ 2$. Thus, $k = 2$. \\
3. \textbf{Calculate product:} $7 \times 2 = 14$.

\vspace{25pt}

Answer: \colorbox{yellow!100}{14}
\end{solution}

\vspace{40pt}

% Question 31
% Category: Practice | Difficulty: Hard | Question Type: Multiple Choice
\question
$f(x) = 2^x$ \\
Which of the following defines a function $g$ whose graph is the result of reflecting the graph of $f$ across the $x$-axis and then shifting it 5 units up?
\begin{tasks}(2)
    \task $g(x) = 2^{-x} + 5$
    \task $g(x) = -2^x + 5$
    \task $g(x) = -2^{x+5}$
    \task $g(x) = 2^x - 5$
\end{tasks}

\begin{solution}
\textbf{Calculation and Logic:} \\
1. \textbf{Reflect across x-axis:} Multiply the whole function by $-1 \rightarrow -f(x) = -2^x$. \\
2. \textbf{Shift up 5:} Add 5 to the result $\rightarrow -2^x + 5$.

\vspace{25pt}

Answer: \colorbox{yellow!100}{B}
\end{solution}

\vspace{40pt}

% Question 32
% Category: Concept | Difficulty: Medium | Question Type: Multiple Choice
\question
\begin{center}
\begin{tikzpicture}[scale=0.6]
    \draw[->] (-1,0) -- (6,0) node[right] {$x$};
    \draw[->] (0,-1) -- (0,5) node[above] {$y$};
    \draw[thick, blue] (1,1) -- (4,4);
    \draw[thick, red, dashed] (2,1) -- (5,4);
    \node[blue] at (1.5,2.5) {$f$};
    \node[red] at (4.5,2.5) {$g$};
\end{tikzpicture}
\end{center}
Which equation describes the relationship between function $f$ and function $g$ shown above?
\begin{tasks}(2)
    \task $g(x) = f(x) + 1$
    \task $g(x) = f(x) - 1$
    \task $g(x) = f(x - 1)$
    \task $g(x) = f(x + 1)$
\end{tasks}

\begin{solution}
\textbf{Conceptual Explanation:} \\
A horizontal movement of the entire graph indicates a shift. Moving to the right is represented by subtracting a constant from the $x$ input.

\vspace{20pt}

\textbf{Calculation and Logic:} \\
The graph of $g$ is the graph of $f$ shifted exactly 1 unit to the right. This is represented by the notation $f(x-1)$.

\vspace{25pt}

Answer: \colorbox{yellow!100}{C}
\end{solution}

\vspace{40pt}

% Question 33
% Category: Exam level | Difficulty: Hard | Question Type: Numeric Entry
\question
The function $f$ is defined by $f(x) = |x-2| + 1$. If $g(x) = f(x+4) - 3$, what is the $x$-coordinate of the vertex (minimum point) of the graph of $g$?

\begin{solution}
\textbf{Calculation and Logic:} \\
1. \textbf{Original vertex of f:} The vertex is $(2, 1)$. \\
2. \textbf{Transformations for g:} $f(x+4)$ shifts the graph 4 units left. $-3$ shifts the graph 3 units down. \\
3. \textbf{Calculate new x:} $2 - 4 = -2$.

\vspace{25pt}

Answer: \colorbox{yellow!100}{-2}
\end{solution}

\vspace{40pt}

% Question 34
% Category: Practice | Difficulty: Medium | Question Type: Multiple Choice
\question
How does the graph of $y = f(x+8)$ compare to the graph of $y = f(x)$?
\begin{tasks}(1)
    \task It is shifted 8 units up.
    \task It is shifted 8 units down.
    \task It is shifted 8 units left.
    \task It is shifted 8 units right.
\end{tasks}

\begin{solution}
\textbf{Conceptual Explanation:} \\
A constant added inside the function parentheses results in a horizontal shift to the left.

\vspace{25pt}

Answer: \colorbox{yellow!100}{C}
\end{solution}

\vspace{40pt}

% Question 35
% Category: Exam level | Difficulty: Hard | Question Type: Numeric Entry
\question
The point $(3, 12)$ is on the graph of $y = f(x)$. What is the $y$-coordinate of the corresponding point on the graph of $y = f(x) - 15$?

\begin{solution}
\textbf{Calculation and Logic:} \\
A transformation of the form $f(x) + k$ only changes the $y$-coordinate. \\
Original $y = 12$. \\
Shift: $-15$. \\
New $y = 12 - 15 = -3$.

\vspace{25pt}

Answer: \colorbox{yellow!100}{-3}
\end{solution}

\vspace{40pt}

% Question 36
% Category: Concept | Difficulty: Medium | Question Type: Multiple Choice
\question
\begin{center}
\begin{tikzpicture}[scale=0.6]
    \draw[->] (-3,0) -- (3,0) node[right] {$x$};
    \draw[->] (0,-3) -- (0,3) node[above] {$y$};
    \draw[thick, blue] (-2,-2) -- (2,2);
    \node at (1,0.5) {$f$};
    \draw[thick, red, dashed] (-2,2) -- (2,-2);
    \node at (1,-1.5) {$g$};
\end{tikzpicture}
\end{center}
Which transformation maps the solid line $f(x) = x$ to the dashed line $g(x) = -x$?
\begin{tasks}(2)
    \task Reflection across the $x$-axis only.
    \task Reflection across the $y$-axis only.
    \task Both A and B are correct for this specific function.
    \task Neither A nor B.
\end{tasks}

\begin{solution}
\textbf{Conceptual Explanation:} \\
For the function $f(x) = x$, a reflection across the $x$-axis gives $-f(x) = -x$, and a reflection across the $y$-axis gives $f(-x) = -x$. For this specific identity line, both transformations result in the same graph.

\vspace{25pt}

Answer: \colorbox{yellow!100}{C}
\end{solution}

\vspace{40pt}

% Question 37
% Category: Exam level | Difficulty: Hard | Question Type: Numeric Entry
\question
The function $f(x) = x^2 + 6x + 9$ is shifted 2 units to the right to create function $g$. What is the value of $g(0)$?

\begin{solution}
\textbf{Calculation and Logic:} \\
1. \textbf{Define g(x):} Shifting 2 units right means $g(x) = f(x-2)$. \\
2. \textbf{Substitute 0 into g:} $g(0) = f(0-2) = f(-2)$. \\
3. \textbf{Evaluate f(-2):} \\
$f(-2) = (-2)^2 + 6(-2) + 9$ \\
$f(-2) = 4 - 12 + 9 = 1$.

\vspace{25pt}

Answer: \colorbox{yellow!100}{1}
\end{solution}

\vspace{40pt}

% Question 38
% Category: Practice | Difficulty: Hard | Question Type: Multiple Choice
\question
If $g(x) = 3f(x-4) + 2$, which of the following is NOT one of the transformations applied to $f$ to get $g$?
\begin{tasks}(1)
    \task A vertical stretch by a factor of 3.
    \task A horizontal shift to the right by 4 units.
    \task A vertical shift up by 2 units.
    \task A reflection across the $x$-axis.
\end{tasks}

\begin{solution}
\textbf{Conceptual Explanation:} \\
Analyze each part of the expression: $3$ is a stretch, $-4$ is a right shift, and $+2$ is an upward shift. No negative multiplier is present, so there is no reflection.

\vspace{25pt}

Answer: \colorbox{yellow!100}{D}
\end{solution}

\vspace{40pt}

% Question 39
% Category: Concept | Difficulty: Medium | Question Type: Multiple Choice
\question
\begin{center}
\begin{tikzpicture}[scale=0.5]
    \draw[->] (-4,0) -- (4,0) node[right] {$x$};
    \draw[->] (0,-1) -- (0,5) node[above] {$y$};
    \draw[domain=-2.2:2.2, smooth, variable=\x, blue, thick] plot ({\x}, {\x*\x});
    \draw[domain=-2.2:2.2, smooth, variable=\x, red, dashed, thick] plot ({\x}, {(\x-3)*(\x-3)});
    \node[blue] at (-1,2) {$f$};
    \node[red] at (3.5,2) {$g$};
\end{tikzpicture}
\end{center}
If $f(x) = x^2$ and $g(x) = (x-c)^2$, what is the value of $c$ based on the graph?
\begin{tasks}(4)
    \task -3
    \task 0
    \task 3
    \task 9
\end{tasks}

\begin{solution}
\textbf{Calculation and Logic:} \\
The vertex of $f$ is at $x = 0$. The vertex of $g$ is at $x = 3$. \\
Since the graph shifted 3 units to the right, the equation must be $(x-3)^2$. Thus, $c = 3$.

\vspace{25pt}

Answer: \colorbox{yellow!100}{C}
\end{solution}

\vspace{40pt}

% Question 40
% Category: Exam level | Difficulty: Hard | Question Type: Numeric Entry
\question
The function $f$ has a minimum value of $-8$. What is the minimum value of the function $g(x) = f(x+10) + 12$?

\begin{solution}
\textbf{Conceptual Explanation:} \\
The horizontal shift $f(x+10)$ does not change the minimum \textit{value} (the $y$-value), it only changes where it occurs. Only the vertical shift $+12$ affects the value.

\vspace{20pt}

\textbf{Calculation and Logic:} \\
Original min $y = -8$. \\
Vertical shift: $+12$. \\
New min $y = -8 + 12 = 4$.

\vspace{25pt}

Answer: \colorbox{yellow!100}{4}
\end{solution}

\vspace{40pt}

% Question 41
% Category: Practice | Difficulty: Medium | Question Type: Multiple Choice
\question
The graph of $y = f(x)$ is reflected across the $y$-axis. Which of the following is the new equation?
\begin{tasks}(4)
    \task $y = -f(x)$
    \task $y = f(-x)$
    \task $y = f(x) - 1$
    \task $y = \frac{1}{f(x)}$
\end{tasks}

\begin{solution}
\textbf{Conceptual Explanation:} \\
Replacing $x$ with $-x$ results in a horizontal reflection over the $y$-axis.

\vspace{25pt}

Answer: \colorbox{yellow!100}{B}
\end{solution}

\vspace{40pt}

% Question 42
% Category: Exam level | Difficulty: Hard | Question Type: Numeric Entry
\question
If the graph of $f(x) = 5x - 2$ is shifted 4 units up and 1 unit left, what is the $y$-intercept of the resulting graph?

\begin{solution}
\textbf{Calculation and Logic:} \\
1. \textbf{Apply Shifts:} $g(x) = f(x+1) + 4$. \\
2. \textbf{Substitute f:} $g(x) = [5(x+1) - 2] + 4$. \\
3. \textbf{Simplify:} $g(x) = 5x + 5 - 2 + 4 = 5x + 7$. \\
4. \textbf{Identify y-intercept:} The constant term is 7.

\vspace{25pt}

Answer: \colorbox{yellow!100}{7}
\end{solution}

\vspace{40pt}

% Question 43
% Category: Concept | Difficulty: Medium | Question Type: Multiple Choice
\question
\begin{center}
\begin{tikzpicture}[scale=0.6]
    \draw[->] (-1,0) -- (5,0) node[right] {$x$};
    \draw[->] (0,-4) -- (0,1) node[above] {$y$};
    \draw[domain=0.5:4.5, smooth, variable=\x, blue, thick] plot ({\x}, {-sqrt(\x)});
    \node at (2,-2.5) {$f$};
\end{tikzpicture}
\end{center}
Which equation could represent the function $f$ shown above if the original function was $y = \sqrt{x}$?
\begin{tasks}(2)
    \task $y = \sqrt{-x}$
    \task $y = -\sqrt{x}$
    \task $y = \sqrt{x} - 4$
    \task $y = \sqrt{x-4}$
\end{tasks}

\begin{solution}
\textbf{Conceptual Explanation:} \\
The graph is below the $x$-axis. This indicates a reflection across the $x$-axis, which is achieved by negating the output of the function.

\vspace{25pt}

Answer: \colorbox{yellow!100}{B}
\end{solution}

\vspace{40pt}

% Question 44
% Category: Exam level | Difficulty: Hard | Question Type: Numeric Entry
\question
The function $g$ is defined by $g(x) = f(x/3)$. If the domain of $f$ is $[0, 6]$, what is the maximum value of the domain of $g$?

\begin{solution}
\textbf{Conceptual Explanation:} \\
A transformation $f(x/c)$ where $c > 1$ results in a horizontal stretch. Every $x$-value in the domain is multiplied by $c$.

\vspace{20pt}

\textbf{Calculation and Logic:} \\
Original domain: $0 \leq x \leq 6$. \\
New domain requirement: $0 \leq x/3 \leq 6$. \\
Multiply by 3: $0 \leq x \leq 18$. \\
The maximum value is 18.

\vspace{25pt}

Answer: \colorbox{yellow!100}{18}
\end{solution}

\vspace{40pt}

% Question 45
% Category: Practice | Difficulty: Medium | Question Type: Multiple Choice
\question
Which transformation of $f(x) = x^2$ would result in a vertex at $(-5, 0)$?
\begin{tasks}(2)
    \task $g(x) = (x-5)^2$
    \task $g(x) = (x+5)^2$
    \task $g(x) = x^2 - 5$
    \task $g(x) = x^2 + 5$
\end{tasks}

\begin{solution}
\textbf{Calculation and Logic:} \\
A vertex at $x = -5$ indicates a shift 5 units to the left. Left shifts are represented by adding the constant inside the parentheses: $(x+5)^2$.

\vspace{25pt}

Answer: \colorbox{yellow!100}{B}
\end{solution}

\vspace{40pt}

% Question 46
% Category: Exam level | Difficulty: Hard | Question Type: Numeric Entry
\question
If $(a, b)$ is a point on $f(x)$, and $(3a, b)$ is a point on $g(x) = f(kx)$, what is the value of $k$?

\begin{solution}
\textbf{Conceptual Explanation:} \\
Multiplying the input by $k$ results in a horizontal scaling. To get from $a$ to $3a$, the graph was stretched horizontally.

\vspace{20pt}

\textbf{Calculation and Logic:} \\
To map $x \rightarrow 3x$, we need $k$ such that when we input $3a$, the function sees $a$. \\
$k(3a) = a \Rightarrow 3k = 1 \Rightarrow k = 1/3$.

\vspace{25pt}

Answer: \colorbox{yellow!100}{1/3}
\end{solution}

\vspace{40pt}

% Question 47
% Category: Concept | Difficulty: Medium | Question Type: Multiple Choice
\question

If $g(x) = \frac{1}{2}f(x)$, how do the $y$-intercepts of $f$ and $g$ compare?
\begin{tasks}(1)
    \task They are the same.
    \task The $y$-intercept of $g$ is half the $y$-intercept of $f$.
    \task The $y$-intercept of $g$ is 2 units higher than $f$.
    \task The $y$-intercept of $g$ is 2 units lower than $f$.
\end{tasks}

\begin{solution}
\textbf{Calculation and Logic:} \\
$y$-intercept occurs at $x=0$. \\
$g(0) = \frac{1}{2}f(0)$. \\
The output is halved.

\vspace{25pt}

Answer: \colorbox{yellow!100}{B}
\end{solution}

\vspace{40pt}

% Question 48
% Category: Exam level | Difficulty: Hard | Question Type: Numeric Entry
\question
The function $f(x) = 10$ is a constant function. What is the value of $f(x+100) + 100$?

\begin{solution}
\textbf{Calculation and Logic:} \\
1. \textbf{Evaluate f(x+100):} Since $f$ is constant, any input results in 10. Thus, $f(x+100) = 10$. \\
2. \textbf{Add 100:} $10 + 100 = 110$.

\vspace{25pt}

Answer: \colorbox{yellow!100}{110}
\end{solution}

\vspace{40pt}

% Question 49
% Category: Practice | Difficulty: Medium | Question Type: Multiple Choice
\question
The graph of $y = f(x)$ is shifted 3 units left and 4 units down. Which of the following is the new equation?
\begin{tasks}(2)
    \task $y = f(x-3) + 4$
    \task $y = f(x+3) - 4$
    \task $y = f(x-3) - 4$
    \task $y = f(x+3) + 4$
\end{tasks}

\begin{solution}
\textbf{Calculation and Logic:} \\
Left 3 $\rightarrow (x+3)$. \\
Down 4 $\rightarrow - 4$.

\vspace{25pt}

Answer: \colorbox{yellow!100}{B}
\end{solution}

\vspace{40pt}

% Question 50
% Category: Exam level | Difficulty: Hard | Question Type: Numeric Entry
\question
The vertex of $f(x) = x^2$ is $(0, 0)$. If $g(x) = -(x+7)^2 - 9$, what is the sum of the coordinates of the vertex of $g$?

\begin{solution}
\textbf{Calculation and Logic:} \\
1. \textbf{Identify Vertex:} The transformation $(x+7)$ means $x = -7$. The transformation $-9$ means $y = -9$. \\
Vertex of g: $(-7, -9)$. \\
2. \textbf{Calculate Sum:} $-7 + (-9) = -16$.

\vspace{25pt}

Answer: \colorbox{yellow!100}{-16}
\end{solution}
% =============================================================
% Question End
% =============================================================

\end{questions}
\begin{center}
\rule{.5\textwidth}{1pt}
\end{center}
\end{document}
